\subsection{I-42. Reset Hessians, sharp block thresholds, and saturated runs}
\label{I-42}

\paragraph{Scope.}
The abstract Hessian examples below are not ROABP annihilators.
The positive polynomial consequence retains its explicit
saturated-run promise.

\paragraph{1. An exact reset-Hessian lemma.}
Let $K$ be a field, let $h_1,\ldots,h_n$ be distinct elements,
and write $v_i=(1,h_i,h_i^2)^{\mathsf T}$. Let $S\subseteq[n]$
be a set of reset positions. Suppose a symmetric zero-diagonal
matrix $H$ has the following form for $i<j$:
\[
 H_{ij}=\begin{cases}
 a_i b_j,&\text{if no element of $S$ lies strictly between $i,j$},\\
 0,&\text{otherwise},
 \end{cases}
 \qquad a_i,b_i\in K.
 \tag{I.147}
\]
Let $t(i)$ be the first reset strictly after $i$, or $n$ if
there is none. Assume
\[
 H\mathbf1=H\mathbf h=H\mathbf h^{\circ2}=0.
 \tag{I.148}
\]
If $t(i)-i\le5$ for every $i<n$, then $H=0$ in every
characteristic. If $\operatorname{char}K\ne2$, the weaker
bound $t(i)-i\le6$ still implies $H=0$.

\paragraph{Proof.}
Suppose $H\ne0$, and let $i$ be its first nonzero row.
All earlier rows and, by symmetry, earlier columns are zero.
Put $t=t(i)$. Some future entry in this row is nonzero,
so $a_i\ne0$. Let $k$ be the first index in $(i,t]$ with
$b_k\ne0$. The three row equations imply
\[
 \sum_{j=k}^t b_jv_j=0.
 \tag{I.149}
\]
Any at most three of the $v_j$ are linearly independent.
Consequently at least four of these $b_j$ are nonzero, and
$t-k\ge3$. Under the first bound $t-i\le5$, this implies
$k-i\le2$.

Since $k\le t-3$, it is not the terminal reset. Every future
entry in its row is still in the same block. Earlier columns
before $i$ are zero, so its row equation is
\[
 b_k\sum_{j=i}^{k-1}a_jv_j+
 a_k\sum_{j=k+1}^t b_jv_j=0.
\]
Use (I.149) and divide by $b_k\ne0$ to obtain
\[
 \sum_{j=i}^{k-1}a_jv_j-a_kv_k=0.
 \tag{I.150}
\]
There are at most three distinct vectors here, so their
independence forces $a_i=0$, a contradiction.

Now suppose $\operatorname{char}K\ne2$ and only $t-i\le6$
is assumed. The same argument applies unless $k=i+3$.
In that remaining case $t=i+6$, and all four coefficients
$b_k,b_{k+1},b_{k+2},b_{k+3}$ are nonzero. Equation (I.150)
is a nonzero relation among four Vandermonde columns; every
coefficient is nonzero, since any three are independent.
In particular $a_k\ne0$, and
\[
 \sum_{j=i}^{k-1}a_jv_j=a_kv_k.
\]
The row at $\ell=k+1<t$ has no contribution beyond this same
block. Substitute the preceding identity and (I.149) into its
equation. This gives
\[
 (2a_kb_\ell-a_\ell b_k)v_k-a_\ell b_\ell v_\ell=0.
 \tag{I.151}
\]
The two columns are independent and $b_\ell\ne0$, so first
$a_\ell=0$ and then $2a_kb_\ell=0$, a contradiction outside
characteristic two. This proves both bounds.

\paragraph{2. Sharpness for abstract Hessians.}
Both thresholds are sharp for matrices of the stated form,
provided the field contains the displayed number of distinct
nodes. In characteristic two, take seven distinct $h_i$.
Choose nonzero vectors $a$ and $b$, supported respectively on
$\{1,2,3,4\}$ and $\{4,5,6,7\}$, with
\[
 \sum_{i=1}^4a_iv_i=0,
 \qquad\sum_{i=4}^7b_iv_i=0.
 \tag{I.152}
\]
Each kernel is one-dimensional and all its displayed entries
are nonzero, by the three-column Vandermonde property. Define
$H_{ij}=a_ib_j$ for $i<j$, impose symmetry, and put $H_{ii}=0$.
The first three rows satisfy (I.148) by the second relation;
the last three do so by the first. The fourth row is
$-2a_4b_4v_4=0$. Thus $H\ne0$ has all three kernels. It is
one reset block of seven positions, with $S=\{1,7\}$, corresponding to five
intermediate invertible positions in the matrix-word setting.

In every characteristic, use eight distinct nodes and take
$a$ supported on $\{1,2,3,4\}$ and $b$ on
$\{5,6,7,8\}$, each annihilating its four-column Vandermonde
matrix. Then
\[
 H=ab^{\mathsf T}+ba^{\mathsf T}
 \tag{I.153}
\]
is a nonzero symmetric zero-diagonal matrix, with $S=\{1,8\}$, of the same
semiseparable form and satisfies (I.148). This proves the
outside-characteristic-two threshold is also exact. These
examples concern an abstract evaluated Hessian; no fixed
polynomial integrability or ROABP realization is asserted.

\paragraph{3. Consequence for an actual rank-preserved width-three word.}
Suppose a nonzero affine multilinear scalar word has at least
one internal rank-two factor and has, at the common
one-pole center $\alpha_i=t+z/(a_i-y)$, all the internal rank
and nonzero active-endpoint properties used in I-41. Assume
its nonzero generic internal interval ranks are two or three,
and that each invertible run has length at most four, or at
most five if the characteristic is not two. Then its nonzero
scalar cannot annihilate the balanced two-pole map.

Indeed, an assumed annihilation gives the same normalized
$I/P$ word, zero first derivatives and passive-zero entries
at singular factors as (I.142)--(I.143). If an endpoint is
followed or preceded by an invertible constant string, transport
its endpoint shear unchanged through that $I_3$ string to the
first or last $P$, exactly as in I-38; this leaves every
constant $I_3$ and $P$ intact. Its Hessian has exactly
the form (I.147), with resets at singular factors and the
endpoint versions of $a_i,b_i$. The distance to the next reset
or final endpoint is at most the invertible-run length plus
one. The three tangent identities give (I.148), so part~1
forces $H=0$. The full state-path argument of I-41 then kills
every possible passive excursion and every output path,
contradicting nonzero scalar output. Thus no additional
Hessian condition remains under these run bounds once the
specified ranks and endpoints have been preserved.

\paragraph{4. A positive saturated-run class beyond weight two.}
Here is one explicit structural application that supplies those
rank hypotheses rather than assumes them. Let $f$ be an
irreducible homogeneous multilinear polynomial, depending on
all its $n$ variables, of width at most three in some order,
and invariant under common translation. Let
$m=n-2\deg f$. Suppose either $m\in\{3,4\}$ in any
characteristic, or $m=5$ in characteristic different from two.
Construct the minimal affine cores from the actual coefficient
support spaces, as in I-40. Assume that every maximal run of
rank-three cores is either isolated or has length exactly $m$.
Then balanced two-pole evaluation hits $f$, in unknown order.
Its finite grid is the same
$(n+1)^2\binom{D+3}{3}$ grid as I-40, with the same field
and extension-input conditions.

To prove this, generic strict internal ranks are at least two
by irreducibility, and the first and last internal factors
are singular by minimality, exactly as in I-41. Isolated
invertible cores give connectors in at most three coordinates,
so I-38's dominance preserves them. Consider a length-$m$
run beginning after cut $j$. Every full-rank core increments
the determinant weight by three, and its entire allowed
budget is $3m$. Hence
\[
 \delta_j=0,\qquad\delta_{j+m}=3m,
 \qquad\delta'_{j+m}=0.
 \tag{I.154}
\]
Its two bounding support spaces are actual three-dimensional
$\mathrm{SL}_2$ submodules. Since $\delta_j=0$, the index
$j$ is even; since $n-m=2D$, the right bounding suffix also
has even length. The adjacent support spaces therefore have
dimension two and odd weights. The argument of I-41,
part~3, excludes all weights other than $\pm1$ and identifies
both bounding three-planes with the actual submodule
$\Gamma^2V\subset V\otimes V$.

Normalize the $m$ full-rank constant cores to the identity.
Each linear coefficient has nonzero entries only on the first
superdiagonal by torus covariance. Their product is still
\[
 M=\begin{pmatrix}1&A&B\\0&1&C\\0&0&1\end{pmatrix},
 \tag{I.155}
\]
where $A,C$ are linear and $B$ quadratic in all $m$ run
variables: a product of three strictly upper triangular
three-by-three matrices is zero. The final left weights are
$(m+2,m,m-2)$ and the right three-plane weights are
$(2,0,-2)$. Semi-invariance of weight $m$ gives the same
nonzero antidiagonal pairing and the same right anchor
as (I.140). Thus the complete connector still has determinant
(I.141), with nonzero constant coefficient $d_2d_3$ on the
square of the left-anchor variable.

Its unique double pole preserves this connector, independently
of how many coordinates occur in $A,B,C$. The other connectors
and univariate pivots survive as above. Therefore every strict
internal interval rank is preserved at the one-pole center.
The minimal endpoint ranks again reduce their active
nonvanishing to pairwise dominance. The maximum run is at
most four in the unrestricted-characteristic cases and at
most five otherwise. Part~3 now proves the asserted scalar
nonvanishing and the claimed finite sampler.

\paragraph{Exact remaining rank gate.}
At weight three, the only missing run length for this
structural theorem is two; at weight four it is two or three.
The reset-Hessian lemma shows that these cases need only a
rank-preservation repair, not an additional bandwidth argument.
The theorem does not claim that the saturated-run promise
survives arbitrary factorization or specialization. It is
stated for an irreducible input and its actual minimal cores.
Products of inputs already known to be hit also remain hit
by the same formal map, since its target is a field, but no
unproved structural promise for their factors is inferred.

The overlapping seven-node characteristic-two example and the
disjoint eight-node example are barriers to extending this
specific Hessian-kernel argument without further information.
They do not disprove a balanced two-pair sampler or a general
polynomial-in-width hitting set. No historical priority claim
is made.
